%% kapa: exam=1 round=37 year=2026 period=2 form=A subject=appraisal_law subject_ko=감정평가관계법규 seq=18 number=18
%% source: Q-Net 감정평가사 기출문제 (cst00309) artlSeq=5258924 file=raw/1008/5258924/2026년 제37회 감정평가사 1차 2교시.pdf
%% conversion: kordoc 4.24.51 -> tools/gen1.py (auto); review=proofread-sonnet; flags=none
\begin{question}{18}
감정평가 및 감정평가사에 관한 법령상 감정평가에 관한 설명으로 옳은 것은?
\begin{choices}
  \item 감정평가법인이 해산한 경우 감정평가서 원본과 그 관련 서류를 해산한 날부터 15일 이내에 국토교통부장관에게 제출해야 한다.
  \item 감정평가서의 적정성 검토업무를 수행할 감정평가사는 5년 이상 감정평가 업무를 수행한 사람으로서 감정평가실적이 100건 이상인 사람이어야 한다.
  \item 국가가 토지의 관리ㆍ매입ㆍ매각ㆍ경매ㆍ재평가 등을 위하여 토지를 감정평가하려는 경우에는 감정평가법인등에 의뢰하지 아니할 수 있다.
  \item 감정평가법인은 의뢰받은 감정평가 결과를 감정평가 정보체계에 등록하여야 하는 경우 감정평가서 발급일부터 30일 이내에 등록해야 한다.
  \item 감정평가 제도를 개선하기 위하여 감정평가서에 대한 표본조사를 실시하는 경우 우선 추출방식에 의할 수 없다.
\end{choices}
\end{question}
\answer{2}
